\section{Diskussion}
I dette afsnit fortolkes projektets resultater med fokus på etiske dilemmaer, metodiske valg og resultaternes validitet i forhold til problemformuleringen.

Et centralt etisk dilemma i projektet har været valget af spillerroller i den digitale oplevelse. Viborg Katedralskoles historie under besættelsen indeholder autentiske beretninger om elevers illegale aktiviteter, herunder smugling af våben, tyveri fra tyske lastbiler og kontakt til modstandsgrupper.\cite{andersen2000viborg} Samtidig beskriver kilderne lærere, der passivt hjalp ved at udlevere nøgler. På baggrund af denne historiske kontekst overvejede projektet initialt et symmetrisk spildesign (fx ”cops and robbers”), hvor brugere kunne påtage sig både modstandsroller og nazistiske roller.

Projektet har bevidst fravalgt at give brugerne mulighed for at spille som nazister. Dette valg er truffet ud fra en etisk vurdering: En gamification af besættelsesmagtens rolle risikerer at trivialisere de historiske rædsler og skabe en upassende identifikation for nuværende elever. Selvom man kunne argumentere for, at en fremstilling, hvor nazisterne taber, også udgør en form for historisk genskrivning, har projektet primært prioriteret at undgå aktiv identifikation med besættelsesmagten. De tildelte ”abilities” i spillet er derfor ikke historisk baserede, men inspireret af eksisterende spilgenrer for at opretholde et funktionelt game loop.

Projektets etiske beslutninger er ikke udførligt dokumenteret i use cases eller user stories. Dette skyldes en bevidst afgrænsning: Projektets hovedfokus har været at demonstrere en sammenhængende digital oplevelse og et velfungerende game loop frem for historisk autenticitet i rollerepræsentationen. Denne prioritering er metodisk forsvarlig inden for projektets tidsramme, men den efterlader et ubehandlet etisk potentiale. Styrken ved denne tilgang er, at projektet leverer et konkret, spilbart produkt. Svagheden er, at brugerens historiske forståelse kan blive overfladisk, og at de etiske implikationer ikke eksplicit formidles i spiloplevelsen.

Set i lyset af problemformuleringens spørgsmål om, ”hvordan bygningen kan fremvises i et historisk perspektiv”, viser projektet, at digitale oplevelser kan formidle historie gennem leg, men at der er et iboende spændingsfelt mellem historisk præcision, etisk ansvar og spilmekanik. Fremtidige projekter bør overveje at involvere målgruppen (nuværende elever) i etiske valg allerede i designfasen, fx gennem workshops om rollerepræsentation.